\section{Analysis}\label{sec:analysis}

% TODO: claim-boundary analysis.  Explain why the result has the scope it has
% and not more.  Subsections in this series typically argue: locality of the
% positive claim, why it does not imply nearby stronger structures, and why
% structural prerequisites (non-faithfulness, non-canonical lifts, etc.) are
% genuine prerequisites rather than separate caveats.

\subsection{Locality of the positive claim}\label{subsec:locality-of-claim}

% TODO: argue that the positive claim is local in the sense the paper
% requires (e.g., layer-local, scope-local, instrument-local).

\subsection{What the result does not imply}\label{subsec:what-not-implied}

% TODO: explain why the result does not on its own license stronger nearby
% structures (event packages, algebraic structures, valuations, etc.).

\subsection{Higher structures are future work}\label{subsec:higher-structures-future-work}

% TODO: identify the structural moves required to reach the disclaimed
% higher structures and frame them as future-work targets, not negative
% results.

\subsection{Structural prerequisite 1}\label{subsec:structural-prerequisite-1}

% TODO: e.g., non-faithful public projection / shadow as a prerequisite for
% the irreducibility statement to be non-trivial.

\subsection{Structural prerequisite 2}\label{subsec:structural-prerequisite-2}

% TODO: e.g., non-canonical selected lift as a complementary prerequisite.
